% !TEX root = ../main.tex
\section{Variational consistency of Kirchhoff plate Galerkin meshfree formulation}\label{sec:formulation}

\subsection{Galerkin weak form of the Kirchhoff plate problem}\label{sec:weakform}

In accordance with the Kirchhoff--Love plate theory \cite{}, the equilibrium equation and the boundary conditions of an elastic plate occupying the domain $\Omega$ with boundary $\Gamma$ can be stated as follows:
\begin{equation}\label{strong}
	\begin{cases}
		\bar D w_{,iijj}=\bar q & \text{in}\;\Omega\\
		w=\bar{w} & \text{on}\;\Gamma_{w}\\
		\theta(w)=w_{,n}=\bar{\theta} & \text{on}\;\Gamma_{\theta}\\
		w=\bar{w} & \text{at}\;c_w\\
		V(w)=-\bar D\left(w_{,iin} +(1-\nu)\left(w_{,ns}\right)_{,s}\right)=\bar{V} & \text{on}\;\Gamma_{V}\\
		M(w)=-\bar D\left(w_{,nn}+\nu w_{,ss}\right)=\bar{M} & \text{on}\;\Gamma_{M}\\
		P(w)=\bar D(1-\nu)\left[w_{,ns}\right]=\bar{P} & \text{at}\;c_P
	\end{cases}
\end{equation}
where the subscript comma of $w_{,i}$ denotes the partial derivative of $w$ with respect to the $i$-th spatial coordinate, $i=1,2$. $\bar D$ denotes the bending rigidity and it can be expressed by Young's modulus $E$, Poisson ratio $\nu$ and plate thickness $t$ as:
\begin{equation}
	\bar D=\frac{Et^{3}}{12(1-\nu^{2})}
\end{equation}
and $\bar q$ the distributed transverse load on the plate.
$\Gamma_{w}$ and $\Gamma_{\theta}$ denote the displacement and slope boundaries, $\bar{w}$ and $\bar{\theta}$ are the corresponding given values.
$\boldsymbol{n}=\{n_i\}$ and $\boldsymbol{s}=\{s_i\}$ represent the boundary outward normal and tangent.
$\bar{V}$ and $\bar{M}$ denote the prescribed equivalent shear force and normal bending moment on the natural boundaries $\Gamma_{V}$ and $\Gamma_{M}$, respectively.
$c_w$ collects the corner points on $\Gamma_{w}$, and $c_P$ collects the corner points where the concentrated corner force $\bar{P}$ is prescribed.
The above boundary segments cover the whole boundary and the corner points of $\Gamma$, denoted by $c$, are shared by them, i.e.
\begin{equation}\label{boundary_relations}
	\Gamma_{w}\cup\Gamma_{V}=\Gamma_{\theta}\cup\Gamma_{M}=\Gamma,\qquad c_w\cup c_P=c
\end{equation}

With the above boundary conditions, the total potential energy functional of the Kirchhoff plate can be expressed as:
\begin{equation}\label{energy_functional}
	\Pi(w) = \int_\Omega \frac{1}{2} \boldsymbol \kappa : \boldsymbol C : \boldsymbol \kappa d\Omega
	- \int_{\Gamma_{V}} w \bar V d\Gamma
	+ \int_{\Gamma_{M}} \theta(w) \bar M d\Gamma
	- w \bar P\vert_{\boldsymbol x \in c_P}
	- \int_\Omega w \bar q d\Omega
\end{equation}
where $\boldsymbol \kappa$ denotes the curvature tensor and $\boldsymbol C$ the elasticity tensor of the plate, their components are expressed as:
\begin{equation}\label{curvature}
	\kappa_{ij} =- w_{,ij},\quad
	C_{ijkl}=\bar D\left(\nu \delta_{ij} \delta_{kl} + \frac{1-\nu}{2}\left(\delta_{ik}\delta_{jl} + \delta_{il}\delta_{jk}\right)\right)
\end{equation}
Using a standard variation operator, i.e. $\delta \Pi (w) = 0$, Eq. \eqref{energy_functional} turns into the variational problem:
\begin{equation}\label{weakform}
	\textrm{Find}\; w \in H_2: \quad a(\delta w, w) = f(\delta w) ,\quad \forall \delta w \in H_2
\end{equation}
with
\begin{equation}\label{bilinear}
	a(\delta w, w) = \int_\Omega \delta \boldsymbol \kappa : \boldsymbol C : \boldsymbol \kappa d\Omega
\end{equation}
\begin{equation}\label{linear}
	f(\delta w) =
	\int_{\Gamma_{V}} \delta w \bar V d\Gamma
	- \int_{\Gamma_{M}} \theta(\delta w) \bar M d\Gamma
	+ \delta w \bar P\vert_{\boldsymbol x \in c_P}
	+ \int_\Omega \delta w \bar q d\Omega
\end{equation}
The above variational problem is well posed provided that the bilinear form $a(\cdot,\cdot)$ is continuous and coercive \cite{}, i.e. there exist two constants $\gamma,\alpha>0$, independent of the nodal distance $h$, such that

\begin{itemize}
\item Continuity
\begin{equation}\label{continuity}
	\left| a(\delta w, w) \right| \le \gamma \Vert \delta w \Vert_{H_2} \Vert w \Vert_{H_2}, \quad \forall \delta w, w \in H_2
\end{equation}

\item Coercivity
\begin{equation}\label{coercivity}
	a(w, w) \ge \alpha \Vert w \Vert_{H_2}^2, \quad \forall w \in H_2
\end{equation}
\end{itemize}
It should be noted that these two properties are the two important conditions for deriving the error estimates of the Galerkin meshfree formulation in what follows.

\subsection{Reproducing kernel meshfree approximation and consistency condition}\label{sec:mls}

In reproducing kernel (RK) meshfree approximation \cite{}, a group of nodes $\{\boldsymbol x_I\}_{I=1}^{n_p}$ is adopted to discretize the problem domain $\Omega$ and the corresponding boundary $\Gamma$, where $n_p$ stands for the total number of meshfree nodes.
The meshfree approximation of the deflection $w(\boldsymbol x)$ $\approx$ $w^{h}(\boldsymbol x)$ is given by:
\begin{equation}\label{approx}
	w^{h}(\boldsymbol x)=\sum_{I=1}^{n_p}\Psi_I(\boldsymbol x) d_I
\end{equation}
with
\begin{equation}\label{shape}
	\Psi_I(\boldsymbol x)=\boldsymbol p^{[p]T}(\boldsymbol x_I-\boldsymbol x) \boldsymbol c(\boldsymbol x) \varphi(\boldsymbol x_I-\boldsymbol x)
\end{equation}
where $d_I$ is the nodal displacement coefficient related to the node $\boldsymbol x_I$, and $\Psi_I(\boldsymbol x)$ is the RK shape function associated with the node $\boldsymbol x_I$ and its compact support domain $\operatorname{supp}(\boldsymbol x_I)$, such that $\Omega\subseteq\bigcup_{I=1}^{n_p}\operatorname{supp}(\boldsymbol x_I)$.
$\boldsymbol p^{[p]}$ in Eq. \eqref{shape} is the $p$-degree monomial basis vector, which is defined as:
\begin{equation}\label{monomial}
	\boldsymbol p^{[p]}(\boldsymbol x)=\left\{1,\ x,\ y,\ \cdots,\ x^{i}y^{j},\ \cdots,\ y^{p}\right\}^{T},\quad 0\le i+j\le p
\end{equation}
It is noted that the compact support $\operatorname{supp}(\boldsymbol x_I)$ comes from the kernel function $\varphi(\boldsymbol x_I-\boldsymbol x)$.
Specifically, the quintic B-spline function \cite{} adopted for constructing the quadratic, cubic and quartic meshfree approximants can be stated as follows:
\begin{equation}\label{kernel}
	\varphi(r)=\frac{1}{5!}\begin{cases}
		(3-3r)^5 - 6(2-3r)^5 + 15(1-3r)^5, & r\le \frac{1}{3} \\
		(3-3r)^5 - 6(2-3r)^5, & \frac{1}{3} < r\le \frac{2}{3} \\
		(3-3r)^5, & \frac{2}{3} < r\le 1 \\
		0, & r> 1 \\
	\end{cases}
\end{equation}
with $r=|\boldsymbol x_I-\boldsymbol x|/s$ being the relative distance between the node $\boldsymbol x_I$ and the evaluated point $\boldsymbol x$, and $s$ is the size of the influence domain.

The unknown vector $\boldsymbol c(\boldsymbol x)$ in Eq. \eqref{shape} is determined by the following $p$-th order consistency conditions:
\begin{equation}\label{consistency}
	\sum_{I=1}^{n_p}\Psi_I(\boldsymbol x) \boldsymbol p^{[p]}(\boldsymbol x_I)=\boldsymbol p^{[p]}(\boldsymbol x)
\end{equation}
or equivalently with a shift form: 
\begin{equation}\label{consistency_shift}
	\sum_{I=1}^{n_p}\Psi_I(\boldsymbol x) \boldsymbol p^{[p]}(\boldsymbol x_I-\boldsymbol x)=\boldsymbol p^{[p]}(\boldsymbol 0)
\end{equation}
Substituting Eq. \eqref{shape} into Eq. \eqref{consistency}, $\boldsymbol c(\boldsymbol x)$ takes the following form:
\begin{equation}\label{coef}
	\boldsymbol c(\boldsymbol x)=\boldsymbol A^{-1}(\boldsymbol x) \boldsymbol p^{[p]}(\boldsymbol 0)
\end{equation}
where $\boldsymbol A(\boldsymbol x)$ is the so-called moment matrix:
\begin{equation}\label{moment}
	\boldsymbol A(\boldsymbol x)=\sum_{I=1}^{n_p}\boldsymbol p^{[p]}(\boldsymbol x_I-\boldsymbol x) \boldsymbol p^{[p]T}(\boldsymbol x_I-\boldsymbol x) \varphi(\boldsymbol x_I-\boldsymbol x)
\end{equation}
Then, combining Eq. \eqref{coef} with Eq. \eqref{moment} yields the following expression of the RK shape functions:
\begin{equation}\label{shape_final}
	\Psi_I(\boldsymbol x)=\boldsymbol p^{[p]T}(\boldsymbol 0) \boldsymbol A^{-1}(\boldsymbol x) \boldsymbol p^{[p]}(\boldsymbol x_I-\boldsymbol x) \varphi(\boldsymbol x_I-\boldsymbol x)
\end{equation}
The meshfree approximation space associated with the RK shape functions is defined as:
\begin{equation}\label{space}
	W^h=\operatorname{span}\{\Psi_I\}_{I=1}^{n_p}
\end{equation}
By virtue of the consistency condition \eqref{consistency}, the RK approximation exactly reproduces any polynomial of degree up to $p$, which leads to the following interpolation error estimate \cite{}:
\begin{equation}\label{interpolation}
	\Vert w - \mathcal I^p w \Vert_{H_k} \le C h^{l-k+1} \vert w \vert_{H_{l+1}}, \quad \forall k \le l+1,\ l \le \min\{p,m\}
\end{equation}
where $\mathcal I^p w\in W^h$ denotes the RK interpolant of $w$, $m$ is the integer such that $w\in H^{m+1}(\Omega)$, and $C$ is a positive constant independent of the nodal distance $h$, which however depends on the normalized support size $s/h$ and the kernel function.

\subsection{Variational consistency and orthogonality}\label{sec:gauss}

Introducing the approximation, like the RK approximation introduced in the previous section, into the variational problem \eqref{weakform} leads to the following Ritz-Galerkin problem:
\begin{equation}\label{discrete}
	\textrm{Find}\; w^h \in W^h: \quad a(\delta w^h, w^h) = f(\delta w^h), \quad \forall \delta w^h \in W^h
\end{equation}
It is noted that the accuracy of the Ritz--Galerkin problem \eqref{discrete} is governed not only by the consistency condition, i.e. the interpolation error estimate \eqref{interpolation} presented in the previous section, but also by another requirement, i.e. the variational consistency \cite{}, which plays a crucial role in the convergence property of Galerkin methods.
A numerical method with approximation space $W^h$ is said to fulfill the $q$-th order variational consistency \cite{} if, while all boundary conditions and external forces are evaluated by arbitrary $q$-th order polynomials, i.e. $v \in P^q$ with $P^q = \operatorname{span}\boldsymbol p^{[q]}$, the following identity holds for a $v \in P^q$:
\begin{equation}\label{variational_consistency}
	a(\delta w^h, v) = b(\delta w^h, v),\quad \forall \delta w^h \in W^h
\end{equation}
where $b$ is an identical functional with $f$ except that the external forces in $f$ are implemented by $v$, as:
\begin{equation}\label{functional_b}
	b(\delta w, v) =
	\int_{\Gamma_{V}} \delta w V(v) d\Gamma
	- \int_{\Gamma_{M}} \theta(\delta w) M(v) d\Gamma
	+ \delta w P(v)\vert_{\boldsymbol x \in c_P}
	+ \int_\Omega \delta w v_{,iijj} d\Omega
\end{equation}

It is noted that the variational consistency of a numerical method is governed by two issues, namely the consistency of the numerical integration and the way in which the essential boundary conditions are imposed.
Most of traditional numerical approximations \cite{}, like finite element methods, have the Kronecker delta property, i.e. $\Psi_I(\boldsymbol x_J) = \delta_{IJ}$.
For these schemes, the essential boundary conditions can be imposed directly, i.e. $\delta w^h=0$ as $\boldsymbol x \in \Gamma_w$ and $\theta(\delta w^h)=0$ as $\boldsymbol x \in \Gamma_{\theta}$, such that Eq. \eqref{variational_consistency} holds for an arbitrary order $q$.
Thereby the boundary integrals in Eq. \eqref{functional_b} can be extended from $\Gamma_{V}$ and $\Gamma_{M}$ to the whole boundary $\Gamma$, and Eq. \eqref{functional_b} can be simplified as:
\begin{equation}\label{functional_b_fem}
	b(\delta w^h, v) =
	\int_{\Gamma} \delta w^h V(v) d\Gamma
	- \int_{\Gamma} \theta(\delta w^h) M(v) d\Gamma
	+ \delta w^h P(v)\vert_{\boldsymbol x \in c_P}
	+ \int_\Omega \delta w^h v_{,iijj} d\Omega
\end{equation}
Substituting Eq. \eqref{functional_b_fem} into Eq. \eqref{variational_consistency}, the variational consistency of the finite element scheme can be verified for an arbitrary order $q$ by the integration by parts.
Since $b(\delta w^h, w)=f(\delta w^h)$ holds for the exact solution $w\in H_2$, the exact solution satisfies the same discrete equation:
\begin{equation}\label{weakform_residual}
	a(\delta w^h, w) = b(\delta w^h, w) = f(\delta w^h), \quad \forall \delta w^h \in W^h
\end{equation}
Meanwhile, the discrete solution $w^h$ satisfies Eq. \eqref{discrete} as well.
Subtracting the two equations from each other gives the orthogonality condition:
\begin{equation}\label{orthogonality}
	a(\delta w^h, w - w^h) = 0, \quad \forall \delta w^h \in W^h
\end{equation}
Thus the orthogonality condition is fully satisfied for FEM.

In contrast to the RK approximation, its shape function can not ensure the Kronecker delta property, $\Psi_I(\boldsymbol x_J) \neq \delta_{IJ}$.
Applying the essential boundary conditions directly like FEM can no longer fulfill the variational consistency of Eq. \eqref{variational_consistency} for an arbitrary order $q$, and thus the orthogonality condition \eqref{orthogonality} is lost.
The variationally-consistent imposition method should be employed herein.
One of the most frequently-used methods is the Nitsche method \cite{}.
The Nitsche method adds consistent terms and stabilized terms in energy functional:
\begin{equation}\label{nitsche_functional}
	\begin{aligned}
		\bar \Pi(w) = \Pi(w) &\underbrace{- \int_{\Gamma_{w}} V(w)(w - \bar w) d\Gamma + \int_{\Gamma_{\theta}} M(w)(\theta(w) - \bar \theta) d\Gamma - P(w)(w - \bar w)\vert_{\boldsymbol x \in c_w}}_{\textrm{Variational terms}} \\
		&\underbrace{+ \frac{\alpha_w}{2} \int_{\Gamma_{w}} (w - \bar w)^2 d\Gamma + \frac{\alpha_\theta}{2} \int_{\Gamma_{\theta}} (\theta(w) - \bar \theta)^2 d\Gamma + \frac{\alpha_c}{2} (w - \bar w)^2\vert_{\boldsymbol x \in c_w}}_{\textrm{Stabilized terms}}
	\end{aligned}
\end{equation}
where $\alpha_w\propto\bar D/h^{3}$, $\alpha_\theta\propto\bar D/h$ and $\alpha_c\propto\bar D/h^{2}$ are the artificial stabilization parameters, which are determined empirically.
With a standard variational operation, the variational problem of Eq. \eqref{nitsche_functional} becomes:
\begin{equation}\label{nitsche_weakform}
	\textrm{Find} \; w^h \in W^h: \quad \bar a(\delta w^h, w^h) = \bar f(\delta w^h), \quad \forall \delta w^h \in W^h
\end{equation}
where $\bar a$ and $\bar f$ denote the bilinear and linear functionals associated with $\bar \Pi(w)$, respectively, i.e.
\begin{equation}\label{nitsche_bilinear}
	\begin{aligned}
		\bar a(\delta w^h, w^h) ={}& a(\delta w^h, w^h)
		- \int_{\Gamma_{w}} \left( V(\delta w^h) w^h + V(w^h) \delta w^h \right) d\Gamma \\
		&+ \int_{\Gamma_{\theta}} \left( M(\delta w^h) \theta(w^h) + M(w^h) \theta(\delta w^h) \right) d\Gamma 
		- \left( P(\delta w^h) w^h + P(w^h) \delta w^h \right) \vert_{\boldsymbol x \in c_w} \\
		&+ \alpha_w \int_{\Gamma_{w}} \delta w^h w^h d\Gamma + \alpha_\theta \int_{\Gamma_{\theta}} \theta(\delta w^h) \theta(w^h) d\Gamma + \alpha_c \delta w^h w^h \vert_{\boldsymbol x \in c_w}
	\end{aligned}
\end{equation}
\begin{equation}\label{nitsche_linear}
	\begin{aligned}
		\bar f(\delta w^h) ={}& f(\delta w^h)
		- \int_{\Gamma_{w}} V(\delta w^h) \bar w d\Gamma
		+ \int_{\Gamma_{\theta}} M(\delta w^h) \bar \theta d\Gamma
		- P(\delta w^h) \bar w \vert_{\boldsymbol x \in c_w} \\
		&+ \alpha_w \int_{\Gamma_{w}} \delta w^h \bar w d\Gamma
		+ \alpha_\theta \int_{\Gamma_{\theta}} \theta(\delta w^h) \bar \theta d\Gamma
		+ \alpha_c \delta w^h \bar w \vert_{\boldsymbol x \in c_w}
	\end{aligned}
\end{equation}

In this context, the variational consistency becomes:
For a $v\in P^q$,
\begin{equation}\label{nitsche_consistency}
	\bar a(\delta w^h, v) = \bar b(\delta w^h, v), \quad \forall \delta w^h \in W^h
\end{equation}
and $\bar b$ denotes the counterpart of $\bar f$ with all the prescribed values implemented by a polynomial $v\in P^q$.
As a result, the consistent term ensures the possibility of enforcement of variational consistency, while the stabilized term recovers the coercivity that is polluted by the consistent term.
Consequently, the bilinear form $\bar a$ is both coercive and continuous, i.e. there exist two constants $\alpha,\gamma>0$, independent of the nodal distance $h$, such that
\begin{equation}\label{nitsche_stability}
	\bar a(w, w)\ge\alpha\Vert w \Vert_{H_2}^2,\quad \forall w\in H_2
\end{equation}
\begin{equation}\label{nitsche_continuity}
	\left|\bar a(\delta w, w)\right|\le\gamma\Vert \delta w \Vert_{H_2}\Vert w \Vert_{H_2},\quad \forall \delta w, w\in H_2
\end{equation}

However, the combination of Gauss quadrature rule and Nitsche method has its problem.
First, the Gauss quadrature rule can not reach the designated accuracy \cite{}, because the overlapping of support size of meshfree shape functions, the shape functions are always the piecewise polynomials and the Gauss quadrature rule is only suitable for polynomial integrals.
As a consequence, the Nitsche formulation can only ensure Eq. \eqref{nitsche_consistency} up to a limited order $q$, which is equal to the order of integration consistency of the quadrature scheme, instead of for an arbitrary order.
For traditional Galerkin meshfree method, the frequently-used Gauss quadrature rules only meet linear variational consistency with $q=1$, while the linear variational consistency is naturally equipped for Kirchhoff palte formualtions.

Second, as can be found in the weak form of Nitsche method, the operation requires the calculation of third derivatives of shape function to evaluate $V$, and the higher-order derivatives of shape functions are complicated and costly.
For the exact solution $w\in H_2$, which is generally not a polynomial, the residual of Eq. \eqref{nitsche_weakform} therefore does not vanish.
To estimate this residual, the $q$-th order Lagrange interpolant $\mathcal T^{q}w$ of $w$ based upon the integration cells is introduced, which satisfies the following estimate \cite{ref50}:
\begin{equation}\label{lagrange_error}
	\Vert w-\mathcal T^{q}w \Vert_{H_{k}}\le C h^{l-k+1}\vert w \vert_{H_{l+1}},\quad \forall k\le l+1,\ l\le\min\{q,m\}
\end{equation}
Since $\mathcal T^{q}w\in P^q$, taking $v=\mathcal T^{q}w$ in Eq. \eqref{nitsche_consistency} gives
\begin{equation}\label{nitsche_consistency_poly}
	\bar a(\delta w^h, \mathcal T^q w) = \bar b(\delta w^h, \mathcal T^q w), \quad \forall \delta w^h \in W^h
\end{equation}
Meanwhile, since $w$ is the exact solution, $\bar b(\delta w^h, w)$ coincides with $\bar f(\delta w^h)$.
Thereby the residual of Eq. \eqref{nitsche_weakform} can be recast as
\begin{equation}\label{defect}
	\bar a(\delta w^h, w) - \bar f(\delta w^h)
	= \bar a(\delta w^h, w) - \bar b(\delta w^h, w) 
	= \bar a(\delta w^h, w - \mathcal T^q w) - \bar b(\delta w^h, w - \mathcal T^q w)
\end{equation}
where the last equality follows from subtracting Eq. \eqref{nitsche_consistency_poly} from the residual.
Following the continuity of $\bar a$ and $\bar b$, the two terms in Eq. \eqref{defect} are bounded by
\begin{equation}\label{defect_bound}
	\left| \bar a(\delta w^h, w - \mathcal T^q w) \right| \le C \Vert \delta w^h \Vert_{H_2} \Vert w - \mathcal T^q w \Vert_{H_2}
\end{equation}
\begin{equation}\label{defect_bound_b}
	\left| \bar b(\delta w^h, w - \mathcal T^q w) \right| \le C \Vert \delta w^h \Vert_{H_2} \Vert w - \mathcal T^q w \Vert_{H_2}
\end{equation}
Substituting Eq. \eqref{lagrange_error} with $k=2$ into Eqs. \eqref{defect_bound} and \eqref{defect_bound_b}, and noting $\bar a(\delta w^h, w^h)=\bar f(\delta w^h)$, the residual \eqref{defect} reduces to the following orthogonality estimate of the Galerkin meshfree formulation:
\begin{equation}\label{orthogonality_estimate}
	\left| \bar a(\delta w^h, w - w^h) \right| \le C h^{l-1} \Vert \delta w^h \Vert_{H_2} \vert w \vert_{H_{l+1}}, \quad \forall l \le \min\{q,m\}
\end{equation}
Eq. \eqref{orthogonality_estimate} measures the loss of the Galerkin orthogonality condition \eqref{orthogonality} induced by the variational inconsistency, and it will serve as a basic ingredient of the error estimates developed in Section \ref{sec:error}.
